\section{Numerical Illustration}
We compare schedules by whole-history projection and protected power. The mechanical setting uses $B=.03313$ N m, $r=4.141\times10^{-4}$ N m/slot, $k=12$, $\delta=.25$ s and four prefix readings. The primary slow-ramp example uses the same specified readout level $\sigma=.01$ N m as the closed-loop illustration; additional noise levels display sensitivity. Physical interpretation and calibration conditions are stated once in Section~\ref{sec:scope}.

\subsection{Slow growth: switching before amplitude escape}\label{sec:slow}
The bias slope is $\kappa=3\times10^{-4}$ N m/s, giving $\eta=7.5\times10^{-5}<r$. The stay target belongs to the finite healthy set through $H=441$, so its minimax power is only $\alpha$ throughout that interval. In the unbounded-level class its information remains zero for every $H$. The symmetric schedule has $\xi=2k\eta/r\approx4.35$ and a first complete block of 36 slots. At $H=300$ (75 s), it uses ten moves and has protected power above $0.999$ at $\alpha=.001$ and $\sigma=.01$, while stay has zero information. This demonstrates useful task contrast under the analytic error radius, before amplitude escape at about 110 s.

\begin{table}[t]\centering\small
\caption{Slow-ramp powers, $\alpha=.001$. $P_i$ is ideal power; $P_g$ uses the same per-side error radius $.002030814$ N m.}\label{tab:slow}
\setlength{\tabcolsep}{2.5pt}
\begin{tabular}{rrrrrrr}\toprule
$H$&$\sigma$&Moves&Stay $P_g$&Fixed-76 $P_g$&Sym. $P_i$&Sym. $P_g$\\\midrule
200&0.01&8&$\alpha^\dagger$&0.065&0.997&0.301\\
300&0.01&10&$\alpha^\dagger$&0.997&$>0.999$&$>0.999$\\
300&0.05&10&$\alpha^\dagger$&0.027&0.301&0.080\\
500&0.05&14&0.0006&0.851&$>0.999$&0.985\\
600&0.10&16&0.004&0.492&0.914&0.749\\
800&0.10&20&0.208&0.997&$>0.999$&$>0.999$\\
1000&0.20&24&0.259&0.941&0.985&0.956\\
\bottomrule\end{tabular}
\par\vspace{2pt}\parbox{.98\columnwidth}{\footnotesize Moves refer to the symmetric schedule. $\dagger$: $G=0$; a constant randomized test rejects with probability $\alpha$. The nonzero-score protection formula is not evaluated. Other rows report the particular protected score, even when its bound falls below $\alpha$.}
\end{table}
The fixed-76 rule is a simple comparator obtained from the local constant-amplitude scale at $a=.1$ N m. It is not refitted at each horizon. At $H=300,\sigma=.01$, it also succeeds, with protected power $.9967$: task switching, rather than the exact growing-hold law alone, explains the main zero-information contrast. The growing-hold rule buys more at intermediate horizons. At $H=500,\sigma=.05$, its lower power is $.9846$ versus $.8514$ for fixed-76. At $H=800,\sigma=.1$, the respective values are $.9996$ and $.9974$; at $H=1000,\sigma=.2$, they are $.9564$ and $.9413$. Thus the protected-power advantage is concentrated in the intermediate tested range; saturation of power does not imply equal asymptotic information efficiency.

Early failures remain relevant. At $H=200,\sigma=.01$ the symmetric ideal power is $.9968$ but protection reduces it to $.3011$. At $H=300,\sigma=.05$ its protected power is only $.0797$. Neither the small pointwise residual nor switching by itself guarantees the requested power. The complete calculation covers eight horizons, three policies, two healthy classes, four noise levels and two test levels. These horizons reach at most 250 s and $.075$ N m, within the analytic envelope of Appendix~\ref{app:embedding}.
\begin{figure}[t]\centering
\includegraphics[width=\columnwidth]{figures/08_slow_power_comparison.pdf}
\caption{Slow-ramp powers at $\sigma=.05$ and $\alpha=.001$. The fixed-76 baseline and the symmetric protected score use identical model-error and health budgets. The ideal symmetric curve separates statistical signal from implementation protection.}\label{fig:slowpower}
\end{figure}

\subsection{Selecting an admission tolerance}
Proposition~\ref{prop:controller} is evaluated at fixed gain 40 over exclusion budgets $k=7,\ldots,24$, with tolerance cap $.01$ N m. The envelope rules out $k=7,8,9$; all budgets $10,\ldots,24$ are feasible. For each budget the symmetric calendar and its whole-history projection are recomputed before data collection. A common 4.5 s initial hold precedes the prefix, so all candidates use an already admitted prefix. This common setup is distinct from the post-onset budget $H$.
\begin{table}[t]\centering\small
\caption{Admission design at $H=500$, $\sigma=.05$, $\alpha=.001$. The radius $e_k$ is the simple exponential bound.}\label{tab:admission}
\setlength{\tabcolsep}{4pt}
\begin{tabular}{rrrrr}\toprule
$k$&Hold (s)&$10^3e_k$ (N m)&Moves&$\mathcal P_k(e_k)$\\\midrule
10&1.00&3.9808&16&0.8843\\
11&1.25&2.1234&16&0.9874\\
12&1.50&2.0309&14&0.9846\\
14&2.00&2.0261&14&0.9736\\
\bottomrule\end{tabular}
\end{table}
The best of the fifteen feasible budgets is $k=11$: the shorter wait compensates for its larger error radius. Increasing $k$ beyond 12 barely changes the floor-dominated error but consumes more observation time. At $H=300,\sigma=.05$, the same best budget raises the bound only from $.0797$ to $.0862$, still far below $.9$. A joint admission choice is a useful finite design operation, not a remedy for an inadequate signal. The baseline results in Table~\ref{tab:slow} remain at $k=12$; the new optimum does not replace them.
\begin{figure}[t]\centering
\includegraphics[width=\columnwidth]{figures/09_admission_tradeoff.pdf}
\caption{The complete feasible admission comparison for $H=500$, $\sigma=.05$, $\alpha=.001$. The optimum is over the stated finite calendar family, not all feedback laws, tolerances or schedules.}\label{fig:admission}
\end{figure}

\subsection{Fast growth: a small residual norm is insufficient}\label{sec:full_disclosure}
A separate faster-ramp example uses $\kappa=.003$ N m/s with the same health and implementation budgets. Table~\ref{tab:guard} includes the low-noise failures and representative higher-noise endpoints. All twelve original ideal first-success records fall below $.9$ after protection; the complete records are given in Supplementary Table S1. Coincident stay/symmetric calendars are grouped here. The low-noise outcomes expose the effect of the coherent penalty $2\varepsilon\|v\|_1/(\sigma\|v\|)$, not merely the ratio $\varepsilon/\sigma$.
\begin{table}[t]\centering\small
\caption{Fast-ramp ideal endpoints: implementation protection. All twelve endpoint records fail the protected $.9$ target.}\label{tab:guard}
\setlength{\tabcolsep}{4pt}
\begin{tabular}{lrrrrr}\toprule
Policy&$\sigma$&$H$&$\alpha$&$P_i$&$P_g$\\\midrule
stay/sym.&0.01&56&$10^{-3}$&0.908&0.079\\
stay/sym.&0.01&70&$10^{-8}$&0.930&0.082\\
sym.&0.10&119&$10^{-3}$&0.905&0.835\\
sym.&0.20&183&$10^{-3}$&0.904&0.861\\
\bottomrule\end{tabular}
\end{table}
These lower guarantees concern a specified linear score under an outer error box; they need not equal its rejection probability on a particular mechanical path. The successful slower-ramp cases and all of these failures use the same analytic error radius, rather than a trajectory-dependent reduced radius.

\subsection{Sharp limits and bounded transient advantage}
For $r=\eta=.1$, $k=2$, $n_0=4$ and $B=\infty$, the exact normalized symmetric-calendar deficits at $H=200,400,800,1600$ are $.5601,.5610,.5623,.5630$, approaching $2\sqrt2/5\approx.5657$. The finite certificate recovers $98.96\%,99.48\%,99.75\%,99.87\%$ of each calendar's exact information. These are specified-calendar projections, not finite-horizon global optima. The constant-amplitude example gives rates $37/192$ and $21/160$ for the marginally optimal hold and the saturation hold. The exact-threshold planned-time approximation gives $124.6,141.1,164.5,189.4$ slots versus exact first successful endpoints $126,142,165,190$ in the separate mathematical example. Gaussian power and preassigned multiple-test levels have finite Monte Carlo implementation checks; the analytical formulas establish the guarantees.
\begin{figure}[t]\centering
\includegraphics[width=\columnwidth]{figures/01_sharp_convergence.pdf}
\caption{Exact information for the symmetric calendar, its finite constructive bound, and the sharp limiting coefficient. The theorem concerns the optimum over deterministic calendars.}\label{fig:sharp}
\end{figure}

For the fast-ramp robot scaling at $\sigma=.2$, the uniform bound in \eqref{eq:explicit_CB} is about $2.97\times10^5$ nats. The dimensional proxy $B^5/(k^2r^2\eta\sigma^2)$ is $53.9$ nats, and the best tested switch-then-stay candidate gains $14.7$ nats over stay. These three values quantify, respectively, a proved calendar-uniform bound, an unproved local-balance scale, and a finite candidate comparison. The first is loose because the proof charges every boundary up to $N_B$ worst-case deficiencies of size $4Ba_j$. The finite-box theorem limits growth of the benefit, not its practical magnitude.

\subsection{Closed-loop reduction checks}
The six-coordinate rigid mechanism was tested with 16 event-aligned extreme paths, 240 random bounded-rate paths and six long histories, in addition to the short illustrative traces. The long histories reach $.75$ N m over 250 s. Their largest admitted residual-vector norm is $1.34\times10^{-4}$ N m, versus the analytic $2.03\times10^{-3}$ N m bound. The latter, not the observed maximum, is used in all protected scores. Noisy-readout calculations and the finite nonlinear stress tests serve different purposes: the former evaluate the stated statistical experiment and the latter check its implementation.

For a 19-reading calendar, the ideal selected-channel information is $32.76$ and the minimax power at $\alpha=10^{-8}$ is $.9935$. A specific admissible path has $.9994$ rejection probability; complete-vector information is $44.27$. Those values demonstrate both the distinction between pathwise and worst-class performance and the additional information available to a complete-vector monitor.
